
\section{\label{sec:intro}Introduction}

Variational quantum algorithms are among the main candidates for testing
optimization routines on near-term quantum devices~\cite{Cerezo2021,Moll2018}.
The quantum approximate optimization algorithm (QAOA) is a standard example: \sam{an initial state is prepared, followed by alternating applications of parametrized cost-phase and mixing unitaries}, 
%a parameterized state is prepared by alternating cost-phase and mixing unitaries, 
and the parameters are optimized classically~\cite{Farhi2014,Blekos2024}.
In its original form, QAOA is naturally suited to unconstrained binary
optimization. Constrained problems are often treated by adding penalty terms to
the cost Hamiltonian, but this embeds the feasible problem into a larger
Hilbert space and may assign probability to infeasible bit strings. The
performance then depends on the relative scale between the objective function
and the penalty terms~\cite{Hadfield2019,Blekos2024}.

Portfolio selection provides a simple and relevant constrained binary
optimization problem. Each binary variable represents whether an asset is
included in the portfolio, and a budget or cardinality constraint fixes the
number of selected assets. Quantum approaches to portfolio optimization have
been studied using QAOA and related variational methods~\cite{Hodson2019,Slate2021, Scursulim2026}.
For fixed-cardinality constraints, a feasible ansatz can be built directly in
the Hamming-weight sector associated with the allowed portfolios. This avoids
sampling infeasible portfolios by construction.

The quantum alternating operator ansatz extends QAOA by allowing mixers adapted
to the feasible set~\cite{Hadfield2019}. A closely related construction is the
quantum walk optimization algorithm (QWOA), in which the mixer is generated by
a quantum walk on a graph whose vertices are feasible configurations 
~\cite{Marsh2019, Slate2021quantumwalkbased}. Different choices of graph define different
feasible mixers. This graph dependence is important: the feasible subspace may
be fixed, while the connectivity used to explore it can still vary.

In this work, we study QWOA mixers for a fixed-cardinality binary portfolio
problem. Our main contribution is to isolate the role of feasible-graph
topology in QWOA for portfolio optimization. In particular, we show that
certain feasible graph topologies, such as Johnson-type graphs, can provide
substantial advantages over the complete feasible graph in the tested
instances. We compare the complete graph on feasible portfolios, the Johnson
graph, and cost-weighted Johnson graphs. The Johnson graph connects two
feasible portfolios when they differ by one asset exchange. In the binary
representation, this corresponds to Hamming distance two, since preserving the
number of selected assets requires flipping one selected and one unselected
asset.

We also clarify the relation between QWOA and the Dicke-state plus
Grover-mixer ansatz. Dicke states are uniform superpositions over
fixed-Hamming-weight strings and are natural initial states for
cardinality-constrained optimization~\cite{Bartschi2019,Cook2019}. Grover
mixers implement phase rotations or reflections around such feasible uniform
states~\cite{Bartschi2020,Bridi2024}. We show that the Dicke-state plus Grover
construction is exactly equivalent, up to a global phase and a parameter
rescaling, to QWOA on the complete graph of feasible portfolios. Thus the
Grover mixer realizes a complete-graph feasible quantum walk.

Our numerical results support a more refined view of feasibility-preserving
variational optimization. In the tested portfolio instances and depths, the
Johnson and lightly weighted Johnson mixers outperform the complete-graph
Grover mixer and a penalty-based QAOA baseline in optimal-sampling probability
and normalized cost gap. This does not establish a general asymptotic
advantage; rather, it provides evidence that the geometry of the feasible graph
is an important design choice in constrained QWOA.

The paper is organized as follows. In Sec.~\ref{sec:dicke-grover}, we define
the feasible-subspace ansatz, introduce QWOA on feasible graphs, define the
weighted Johnson extension, and prove the equivalence between complete-graph
QWOA and the Dicke-state plus Grover mixer. In Sec.~\ref{sec:simulation}, we
describe the portfolio data and simulation protocol. In Sec.~\ref{sec:results},
we compare the feasible mixers and the penalty-based QAOA formulation.
Section~\ref{sec:conclusion} gives our conclusions.
